% !TeX root = ../main.tex

\chapter{历史分段与收益持续性}
\section{初始期、成熟期与追加区间}
\begin{table}[htbp]
\centering
\caption{不同历史区间的年化收益（\%）}
\renewcommand{\arraystretch}{1.18}
\small
\setlength{\tabcolsep}{3pt}
\begin{tabularx}{\textwidth}{@{}p{3.15cm}*{6}{C}@{}}
\toprule
区间 & E & P1 & P2 & P3 & P4 & 上证 \\
\midrule
2003-04至2008-03 & 70.52 & 70.52 & 70.52 & 45.25 & 45.25 & 18.11 \\
2008-04至2015-06 & 20.53 & 24.76 & 27.14 & 21.38 & 21.86 & 3.27 \\
2015-07至2024-12 & 9.95 & 15.10 & 12.19 & 13.26 & 10.65 & -2.62 \\
2025-01至2026-09 & 36.03 & 32.53 & 11.18 & 29.07 & 11.48 & 8.12 \\
\bottomrule
\end{tabularx}
\end{table}

初始期60个月，E、P1、P2均年化70.52\%，反映三者完全相同的早期选股和特定市场路径。其高复利对全期财富积累具有较大影响，不能把全期约30\%的年化描述为策略进入成熟市场后一直保持的增长速度。风控路径在同段降为45.25\%，也显示缩仓会放弃部分上涨收益。

成熟期前段P2年化27.14\%，高于P1的24.76\%；后续2015年7月至2024年12月则P1为15.10\%、P2为12.19\%，排序逆转。追加21个月中，E年化36.03\%，高于P1的32.53\%、P3的29.07\%、P4的11.48\%和P2的11.18\%。该区间较短，年化换算对短期路径敏感，但已足以反驳“四条优化在每一时期都优于基础延展”的表述。

分段回撤在各段初始净值1上重新计算，不能直接替代跨阶段的全期回撤。历史分段、参数选择和研究者已见数据的关系也限制预测解释，本表不承担严格样本外验证功能。

\section{年度表现与滚动窗口}
\begin{figure}[htbp]
\centering
\includegraphics[width=\linewidth]{fig08}
\caption{年度复利收益的跨方案比较}
\par\smallskip
\begin{minipage}{0.97\linewidth}\small 注：数字为实际年度复利百分比，色阶仅为便于阅读截于$\pm$60\%；2003年与2026年各含9个月。资料来源：CSMAR，作者计算。\end{minipage}
\end{figure}

年度热图显示，同一时期的市场环境对各路径有共同影响。风控路径并不逐年增加收益，选股路径的相对排序也随年份变化。完整年度收益详列附录A，避免只呈现有利年份。2026年结果必须与“前9个月”的统计长度一起解读。

\begin{figure}[htbp]
\centering
\includegraphics[width=\linewidth]{fig09}
\caption{滚动五年年化收益}
\par\smallskip
\begin{minipage}{0.97\linewidth}\small 注：每个点使用此前60个月；窗口相互重叠，不能作为彼此独立的观察或胜率检验样本。资料来源：CSMAR，作者计算。\end{minipage}
\end{figure}

滚动曲线用于显示收益水平随历史终点变化，而不是重复筛选最优参数。高全期CAGR与较低后段收益可以同时成立；以固定全期均值替代各阶段经验，会掩盖早期复利与后期回报的差别。

\section{初始高收益与后续延展应当怎样连接}
初始期60个月的高收益来自三条未控规则完全相同的选股记录，因此不能作为成熟期优化的证据。它会共同抬升E、P1和P2的完整区间年化，使一个成熟期表现一般的路径仍呈现较高全期数值。为此，成熟期合计还应独立阅读：E、P1、P2年化分别为16.30\%、20.39\%、17.73\%，均低于各自包含初始期的完整年化。

这不是把早期收益删除，而是区分完整投资路径与规则改变后阶段。2007年12月产生的初始规则名单持续影响2008年1至3月，所以分界必须按实际持有月份放在2008年4月，不能仅按自然年2008年1月切开。此处的月份连接与程序的“季度末信号、后三个月持有”一致。

2015年7月至2024年12月，P1年化15.10\%高于P2的12.19\%；追加21个月中，E以36.03\%领先，P1为32.53\%、P2仅11.18\%。原区间与延展区间的排名逆转有实际数据依据，因此本研究同时保留两条路径，不按某一个终点重新决定唯一正确策略。追加区间亦是已知历史的一部分，不能被包装为事前冻结后的独立预测验证。

\section{年度优势与回撤恢复的补充观察}
附录D进一步列出完整自然年度的正收益数、领先基础延展的年度数以及最差年度；只纳入2004—2025年各12个月的完整年度。2003年和2026年不混入年度计数，避免长度差异。年度次数只是描述，对照中的年份也共享市场环境，不能解释为未来获胜概率或独立试验。

同时，最大回撤幅度之外还应关注从峰值到谷值、再到首次恢复峰值的时间。一条路径即使最大损失较小，也可能较长时间未创新高。附录D以每月末净值计算相关日期和月数，初始净值纳入峰值搜索；这比单一回撤百分数多提供时间维度，但仍无法观察月内低点或投资者实际申赎路径。
